\documentclass{beamer}

%lembre-se de configurar o pdflatex com:
%pdflatex -shell-escape -synctex=1 -interaction=nonstopmode  -output-directory=build  %.tex

\usepackage{aula}

\title{Processamento paralelo básico}
\date{\today}

\begin{document}
    
\input{capa}

\begin{frame}[fragile]{Threads vs Processos (Python)}
    
    \begin{columns}[t]
        
        \column{0.48\textwidth}
        \textbf{Threads}
        
        \begin{itemize}
            \item Compartilham memória
            \item Baixo custo de criação
            \item Executam dentro do mesmo processo
            \item Limitadas pelo \textbf{GIL}
        \end{itemize}
        
        \textbf{Uso típico:}
        \begin{itemize}
            \item I/O-bound (download, arquivos, APIs)
        \end{itemize}
        
        \column{0.48\textwidth}
        \textbf{Processos}
        
        \begin{itemize}
            \item Memória isolada
            \item Mais pesados (fork/spawn)
            \item Executam em paralelo real
            \item Não sofrem com o GIL\\
            
        \end{itemize}
        
        \textbf{Uso típico:}
        \begin{itemize}
            \item CPU-bound (cálculos, raster, ML)
        \end{itemize}
        
    \end{columns}
    
    \vspace{0.3cm}
    
    \textbf{Resumo}
    
    \begin{itemize}
        \item Threads → rápidas, mas limitadas pelo GIL
        \item Processos → paralelismo real, porém mais custo
    \end{itemize}
    
\end{frame}

\begin{frame}[fragile]{concurrent.futures — Conceito}
    
    O módulo \texttt{concurrent.futures} fornece uma API de alto nível para \textbf{execução paralela}.
    
    \vspace{0.4cm}
    
    Principais executores:
    
    \begin{itemize}
        \item \textbf{ThreadPoolExecutor} — múltiplas threads
        \item \textbf{ProcessPoolExecutor} — múltiplos processos
    \end{itemize}
    
    \vspace{0.4cm}
    
    \textbf{Quando usar:}
    
    \begin{itemize}
        \item I/O-bound → threads (download, leitura de arquivos)
        \item CPU-bound → processos (cálculos pesados)
    \end{itemize}
    
\end{frame}

\begin{frame}[fragile]{Executando tarefas com ThreadPoolExecutor}
    
    Exemplo: executar várias tarefas em paralelo.
    
    \begin{minted}{python}
from concurrent.futures import ThreadPoolExecutor
import time

def tarefa(x):
    time.sleep(1)
    return x * 2

with ThreadPoolExecutor(max_workers=4) as executor:
    resultados = list(executor.map(tarefa, range(5)))

    print(resultados)
    \end{minted}
    
    \textbf{map():}
    \begin{itemize}
        \item Executa a função para cada elemento
        \item Retorna os resultados na mesma ordem
    \end{itemize}
    
\end{frame}

\begin{frame}[fragile]{Futures e submit()}
    
    \texttt{submit()} permite controle mais detalhado.
    
    \begin{minted}{python}
from concurrent.futures import ThreadPoolExecutor

def tarefa(x):
    return x * x

with ThreadPoolExecutor() as executor:
    futures = [executor.submit(tarefa, i) for i in range(5)]

for f in futures:
    print(f.result())
    \end{minted}
    
    \vspace{0.3cm}
    
    \textbf{Future:}
    
    \begin{itemize}
        \item Representa uma tarefa assíncrona
        \item Métodos:
        \begin{itemize}
            \item \texttt{result()}
            \item \texttt{done()}
            \item \texttt{exception()}
        \end{itemize}
    \end{itemize}
    
\end{frame}

\begin{frame}[fragile]{ProcessPoolExecutor (CPU-bound)}
    
    Para tarefas pesadas, use múltiplos processos.
    
    \begin{minted}{python}
from concurrent.futures import ProcessPoolExecutor

def pesado(x):
    s = 0
    for i in range(10**7):
    s += i * x
    return s

with ProcessPoolExecutor() as executor:
    resultados = list(executor.map(pesado, range(4)))

    print(resultados)
    \end{minted}
    
    \textbf{Importante:}
    
    \begin{itemize}
        \item Evita o GIL (Global Interpreter Lock)
        \item Cada processo roda em um núcleo diferente
    \end{itemize}
    
\end{frame}

\begin{frame}[fragile]{Exemplo: Download paralelo com ThreadPoolExecutor}
    
    \begin{minted}{python}
from concurrent.futures import ThreadPoolExecutor
import requests
from pathlib import Path

urls = ["https://example.com/file1.zip", "https://example.com/file2.zip"]

output_dir = Path("downloads")
output_dir.mkdir(exist_ok=True)

def download(url):
    nome = url.split("/")[-1]
    caminho = output_dir / nome

    r = requests.get(url)
    if r.status_code == 200:
        with open(caminho, "wb") as f:
            f.write(r.content)
            print(f"{nome} OK")
    else:
            print(f"Erro: {nome}")

with ThreadPoolExecutor(max_workers=4) as executor:
    executor.map(download, urls)
    \end{minted}
    
\end{frame}

\begin{frame}[fragile]{asyncio — Conceito}
    
    \textbf{asyncio} permite programação \textbf{assíncrona} em Python usando \textbf{corrotinas}.
    
    \vspace{0.4cm}
    
    \textbf{Ideia central}
    
    \begin{itemize}
        \item Executar tarefas sem bloquear o programa
        \item Ideal para operações \textbf{I/O-bound}
        \item Baseado em \textbf{event loop}
    \end{itemize}
    
    \vspace{0.4cm}
    
    \textbf{Principais conceitos}
    
    \begin{itemize}
        \item \texttt{async def} → define corrotina
        \item \texttt{await} → pausa execução até resultado
        \item \texttt{event loop} → gerencia execução
    \end{itemize}
    
\end{frame}

\begin{frame}[fragile]{Exemplo básico}
    
    \begin{minted}{python}
import asyncio

async def tarefa():
    print("Início")
    await asyncio.sleep(1)
    print("Fim")

asyncio.run(tarefa())
    \end{minted}
    
    \vspace{0.3cm}
    
    \textbf{Observações}
    
    \begin{itemize}
        \item \texttt{sleep()} não bloqueia o programa
        \item Outras tarefas podem rodar durante o \texttt{await}
    \end{itemize}
    
\end{frame}

\begin{frame}[fragile]{Executando múltiplas tarefas}
    
    \begin{minted}{python}
import asyncio

async def tarefa(nome, tempo):
    print(f"Início {nome}")
    await asyncio.sleep(tempo)
    print(f"Fim {nome}")

async def main():
    await asyncio.gather(
        tarefa("A", 1),
        tarefa("B", 2),
        tarefa("C", 1)
    )

asyncio.run(main())
    \end{minted}
    
    \vspace{0.3cm}
    
    \textbf{gather():}
    \begin{itemize}
        \item Executa várias corrotinas "em paralelo"
        \item Retorna quando todas terminam
    \end{itemize}
    
\end{frame}

\begin{frame}[fragile]{Exemplo prático: múltiplas requisições}
    
    \begin{minted}{python}
import asyncio
import aiohttp

urls = [
    "https://example.com/1",
    "https://example.com/2",
    "https://example.com/3"
]

async def fetch(session, url):
    async with session.get(url) as resp:
        return await resp.text()

async def main():
    async with aiohttp.ClientSession() as session:
        tarefas = [fetch(session, u) for u in urls]
        resultados = await asyncio.gather(*tarefas)

print(len(resultados))

asyncio.run(main())
    \end{minted}
    
    \textbf{Vantagem:} alta eficiência em requisições HTTP simultâneas.
    
\end{frame}

\begin{frame}[fragile]{Threads vs asyncio (comparação direta)}
    
    \begin{columns}[t]
        
        \column{0.48\textwidth}
        \textbf{Threads}
        
        \begin{itemize}
            \item Paralelismo via \textbf{múltiplas threads}
            \item Uso do SO (context switch)
            \item Mais simples de usar
            \item Pode ter \textbf{race conditions}
            \item Limitado pelo GIL (CPU)
        \end{itemize}
        
        \begin{minted}{python}
from concurrent.futures import ThreadPoolExecutor

def tarefa(x):
    return x * 2

with ThreadPoolExecutor() as ex:
    resultados = list(ex.map(tarefa, range(5)))
        \end{minted}
        
        \textbf{Modelo:}
        \begin{itemize}
            \item Preemptivo (SO decide troca)
        \end{itemize}
        
        \column{0.48\textwidth}
        \textbf{asyncio}
        
        \begin{itemize}
            \item Concorrência via \textbf{corrotinas}
            \item Single-thread (event loop)
            \item Mais eficiente para I/O massivo
            \item Sem race condition por padrão
            \item Requer \texttt{async/await}
        \end{itemize}
        
        \begin{minted}{python}
import asyncio

async def tarefa(x):
    return x * 2

async def main():
    res = await asyncio.gather(
        *(tarefa(i) for i in range(5))
    )

asyncio.run(main())
        \end{minted}
        
        \textbf{Modelo:}
        \begin{itemize}
            \item Cooperativo (await libera execução)
        \end{itemize}
        
    \end{columns}
    
    \vspace{0.3cm}
    
    \textbf{Resumo prático}
    
    \begin{itemize}
        \item Poucas tarefas I/O → tanto faz
        \item Muitas conexões (HTTP, APIs) → \textbf{asyncio}
        \item Código legado ou simples → \textbf{threads}
    \end{itemize}
    
\end{frame}

\end{document}
